\hdhead{40}{多边形的认识——信息窗1}
\hdsection{1}{146}{连一连。}
\begin{scope}[scale=.6666667,draw=black,text=black]\draw[hdoutline] (110.000,243.000)--(146.000,188.000)--(207.000,243.000)--cycle;\end{scope}
\begin{scope}[scale=.6666667,draw=black,text=black]\draw[hdoutline] (232.000,191.000)--(265.000,240.000)--(336.000,240.000)--cycle;\end{scope}
\begin{scope}[scale=.6666667,draw=black,text=black]\draw[hdoutline] (375.000,188.000)--(444.000,188.000)--(375.000,244.000)--cycle;\end{scope}
\begin{scope}[scale=.6666667,draw=black,text=black]\draw[hdoutline] (473.000,190.000)--(473.000,241.000)--(580.000,241.000)--cycle;\end{scope}
\begin{scope}[scale=.6666667,draw=black,text=black]\draw[hdoutline] (610.000,242.000)--(635.000,188.000)--(660.000,242.000)--cycle;\end{scope}
\begin{scope}[scale=.6666667,draw=black,text=black]\draw[hdoutline] (689.000,235.000)--(753.000,194.000)--(809.000,235.000)--cycle;\end{scope}
\begin{scope}[scale=.6666667,draw=black,text=black]\draw[line width=.45pt] (383.000,188.000)--(383.000,196.000)--(375.000,196.000);\end{scope}
\begin{scope}[scale=.6666667,draw=black,text=black]\draw[line width=.45pt] (473.000,233.000)--(481.000,233.000)--(481.000,241.000);\end{scope}
\q{118}{344}{180}{13}{钝角三角形}
\q{426}{344}{185}{13}{直角三角形}
\q{685}{344}{180}{13}{锐角三角形}
\hdsection{2}{402}{分别画出下面三角形底边上的高。}
\begin{scope}[scale=.6666667,draw=black,text=black]\draw[hdoutline] (164.000,446.000)--(233.000,538.000)--(352.000,532.000)--cycle;\end{scope}
\begin{scope}[scale=.6666667,draw=black,text=black]\draw[hdoutline] (451.000,443.000)--(451.000,535.000)--(543.000,534.000)--cycle;\end{scope}
\begin{scope}[scale=.6666667,draw=black,text=black]\draw[hdoutline] (642.000,516.000)--(699.000,444.000)--(755.000,517.000)--cycle;\end{scope}
\begin{scope}[scale=.6666667,draw=black,text=black]\node[inner sep=0pt,font=\normalsize,text=black] at (273.000,474.000) {底};\end{scope}
\begin{scope}[scale=.6666667,draw=black,text=black]\node[inner sep=0pt,font=\normalsize,text=black] at (514.000,474.000) {底};\end{scope}
\begin{scope}[scale=.6666667,draw=black,text=black]\node[inner sep=0pt,font=\normalsize,text=black] at (699.000,529.000) {底};\end{scope}
\hdsection{3}{574}{画一画，标一标。}

\q{99}{660}{735}{12}{和边。}
\hdsection{4}{842}{按要求画一画。}
\q{101}{886}{231}{12}{① 既是直角三角形，}
\q{101}{926}{231}{12}{　又是等腰三角形。}
\q{355}{886}{233}{12}{② 既是钝角三角形，}
\q{355}{926}{233}{12}{　又是等腰三角形。}
\q{588}{886}{249}{12}{③ 既是锐角三角形，}
\q{588}{926}{249}{12}{　又是等腰三角形。}
\begin{scope}[scale=.6666667,draw=black,text=black]\draw[line width=.5pt] (130,967) rectangle (317,1154);\end{scope}
\begin{scope}[scale=.6666667,draw=black,text=black]\foreach \i in {0,...,9}{\foreach \j in {0,...,9}{\fill[black] ({143.5+\i*17.5},{979.5+\j*17.5}) circle (1.15);}}\end{scope}
\begin{scope}[scale=.6666667,draw=black,text=black]\draw[line width=.5pt] (385,967) rectangle (573,1154);\end{scope}
\begin{scope}[scale=.6666667,draw=black,text=black]\foreach \i in {0,...,9}{\foreach \j in {0,...,9}{\fill[black] ({399.5+\i*17.5},{979.5+\j*17.5}) circle (1.15);}}\end{scope}
\begin{scope}[scale=.6666667,draw=black,text=black]\draw[line width=.5pt] (618,967) rectangle (812,1154);\end{scope}
\begin{scope}[scale=.6666667,draw=black,text=black]\foreach \i in {0,...,9}{\foreach \j in {0,...,9}{\fill[black] ({632.5+\i*17.5},{979.5+\j*17.5}) circle (1.15);}}\end{scope}
